\section{Best Time to Buy and Sell Stock I}
\label{sec:dp-stock-1}

\problemstatement{Given an array $\textit{price}$ where
$\textit{price}[i]$ is a stock's price on day $i$, find the maximum
profit from \textbf{at most one} buy followed by one sell (the buy must
happen before the sell).}

\begin{patternbox}
\emph{``At most one transaction''} is the simplest member of
Section~\ref{sec:dp-stocks-intro}'s holding-state family: selling is a
\textbf{terminal} action here, since no second buy is ever allowed
afterward.
\end{patternbox}

\subsection*{Deriving the recurrence}

Using the base template, the only change is the \textbf{sell} clause:
selling on day $i$ ends all activity for good, contributing exactly
$\textit{price}[i]$ and nothing more (not $f(i{+}1,0)$, which would
incorrectly allow a second buy):
\[
  f(i, \textit{holding}) =
  \begin{cases}
    \max\big(-\textit{price}[i] + f(i{+}1, 1),\ f(i{+}1, 0)\big) &
    \text{not holding}, \\
    \max\big(\textit{price}[i],\ f(i{+}1, 1)\big) & \text{holding.}
  \end{cases}
\]
Base case: $f(n, \cdot) = 0$.

\subsection*{Approach 1: Recursion}

\begin{lstlisting}
#include <bits/stdc++.h>
using namespace std;

int stock1Recursive(int i, int holding, vector<int>& price) {
    int n = price.size();
    if (i == n) return 0;

    if (holding) {
        return max(price[i], stock1Recursive(i + 1, 1, price));
    }
    return max(-price[i] + stock1Recursive(i + 1, 1, price),
               stock1Recursive(i + 1, 0, price));
}

int main() {
    vector<int> price = {7, 1, 5, 3, 6, 4};
    cout << stock1Recursive(0, 0, price) << endl; // 5
    return 0;
}
\end{lstlisting}

\begin{complexitybox}[Approach 1 Complexity]
\textbf{Time.} $O(2^n)$.
\textbf{Space.} $O(n)$ recursion stack.
\end{complexitybox}

\subsection*{Approach 2: Memoization}

\begin{lstlisting}
#include <bits/stdc++.h>
using namespace std;

int stock1Memo(int i, int holding, vector<int>& price, vector<vector<int>>& memo) {
    int n = price.size();
    if (i == n) return 0;
    if (memo[i][holding] != -1) return memo[i][holding];

    int result;
    if (holding) {
        result = max(price[i], stock1Memo(i + 1, 1, price, memo));
    } else {
        result = max(-price[i] + stock1Memo(i + 1, 1, price, memo),
                      stock1Memo(i + 1, 0, price, memo));
    }
    return memo[i][holding] = result;
}

int main() {
    vector<int> price = {7, 1, 5, 3, 6, 4};
    int n = price.size();
    vector<vector<int>> memo(n, vector<int>(2, -1));
    cout << stock1Memo(0, 0, price, memo) << endl; // 5
    return 0;
}
\end{lstlisting}

\begin{complexitybox}[Approach 2 Complexity]
\textbf{Time.} $O(n)$ (only $2$ values of \textit{holding}).
\textbf{Space.} $O(n)$ memo $+$ $O(n)$ recursion stack.
\end{complexitybox}

\subsection*{Approach 3: Tabulation}

\begin{lstlisting}
#include <bits/stdc++.h>
using namespace std;

int stock1Tabulation(vector<int>& price) {
    int n = price.size();
    vector<vector<int>> dp(n + 1, vector<int>(2, 0));

    for (int i = n - 1; i >= 0; i--) {
        dp[i][1] = max(price[i], dp[i + 1][1]);
        dp[i][0] = max(-price[i] + dp[i + 1][1], dp[i + 1][0]);
    }
    return dp[0][0];
}

int main() {
    vector<int> price = {7, 1, 5, 3, 6, 4};
    cout << stock1Tabulation(price) << endl; // 5
    return 0;
}
\end{lstlisting}

\subsection*{Dry run of the tabulation table}

\dryrun{Take $\textit{price} = [7,1,5,3,6,4]$.}

Filling backward from day $5$: $\textit{dp}[5][1]=\max(4,0)=4$,
$\textit{dp}[5][0]=\max(-4+0,0)=0$ (buying on the last day with no
chance to sell is never better than doing nothing, so this clause
$-4+\textit{dp}[6][1]=-4+0=-4$ loses to $0$). Continuing backward, by
day $1$: $\textit{dp}[1][1]=\max(1,\ldots)$ tracks the best sell-from-here
value, and $\textit{dp}[1][0]=\max(-1+\textit{dp}[2][1], \ldots)$
captures buying at price $1$ then later selling at the best subsequent
holding value, $6$, giving profit $5$. The final answer
$\textit{dp}[0][0]=5$ matches buying at price $1$ (day $1$) and selling
at price $6$ (day $4$).

\begin{complexitybox}[Approach 3 Complexity]
\textbf{Time.} $O(n)$.
\textbf{Space.} $O(n)$ table, no recursion stack.
\end{complexitybox}

\subsection*{Approach 4: Space Optimization}

\begin{lstlisting}
#include <bits/stdc++.h>
using namespace std;

int stock1SpaceOptimized(vector<int>& price) {
    int n = price.size();
    int nextHold = 0, nextNotHold = 0;

    for (int i = n - 1; i >= 0; i--) {
        int curHold = max(price[i], nextHold);
        int curNotHold = max(-price[i] + nextHold, nextNotHold);
        nextHold = curHold;
        nextNotHold = curNotHold;
    }
    return nextNotHold;
}

int main() {
    vector<int> price = {7, 1, 5, 3, 6, 4};
    cout << stock1SpaceOptimized(price) << endl; // 5
    return 0;
}
\end{lstlisting}

\begin{complexitybox}[Approach 4 Complexity]
\textbf{Time.} $O(n)$.
\textbf{Space.} $O(1)$ -- just two scalars per day, since the
\textit{holding} dimension already has only $2$ values.
\end{complexitybox}

\begin{ideabox}[The Even Simpler Greedy Alternative]
Because only one transaction is allowed, this specific variant also
admits an $O(n)$, $O(1)$ greedy one-pass solution: track the minimum
price seen so far while scanning left to right, and at each day compute
$\textit{price}[i] - \textit{minSoFar}$, keeping the running maximum.
This works only because there is no transaction-count or cooldown
complexity to track -- the DP formulation above is what generalizes to
every other variant in this chapter, which is why it is presented first.
\end{ideabox}

\begin{takeawaybox}
The holding-state DP template, even when a simpler greedy alternative
exists for this specific variant (as it does here), is worth learning
first specifically because it is the version that survives every
additional constraint the rest of this chapter introduces.
\end{takeawaybox}
